\documentclass{article}
\usepackage{amsmath, amssymb}
\usepackage{graphicx}
\usepackage{float}
\begin{document}
	\section{SCASTIE}
	Here I used SCASTIE to run the provided program.
	
	\begin{figure}[H]
		\centering
		\includegraphics[width=0.9\textwidth]{scastie_image.png}
		\caption{SCASTIE program output for initial program}
	\end{figure}
	
	Getting SCASTIE running was no real issue for me. It was easy to find and made writing code very easy.
	
		\newpage
	
	\begin{figure}[H]
		\centering
		\includegraphics[width=0.9\textwidth]{scastie_image_2.png}
		\caption{SCASTIE program output for updated sum}
	\end{figure}
	
	This program calculates the sum of integers from 21 through 137. The correctness of this program can be verified by looking at the closed-form formula for the sum of consecutive integers.
	
	\[
	\sum_{n=21}^{137}n = \frac{137(137+1)}{2} - \frac{20(20+1)}{2} 
	\] 
	
	which comes out to:
	\[
	= \frac{18906}{2} - \frac{420}{2}
	= 9453 - 210
	= 9243
	\]
	
	\newpage
	
	\begin{figure}[H]
		\centering
		\includegraphics[width=0.9\textwidth]{zed_image.png}
		\caption{Scala program on local machine}
	\end{figure}
	
	\begin{figure}[H]
		\centering
		\includegraphics[width=0.9\textwidth]{output.png}
		\caption{Output for local program}
	\end{figure}
	
	Installing Scala wasn't difficult, although I did run into a small issue after getting it installed initially. Scala and Coursier had to be manually added to my PATH so that their commands would be recognized by the terminal.
	
\end{document}